\chapter{Usage of Chatbots and Other Tools}
\label{ch:ai-tools}

AI-based assistants were used as supplementary software-development tools,
consistent with the PPP guideline. Their roles included:
\begin{itemize}
  \item explaining library error messages and API behavior;
  \item proposing localized code changes and diagnostic tests;
  \item reviewing mathematical notation and tensor-index conventions;
  \item suggesting report organization and language revisions; and
  \item identifying inconsistencies requiring independent investigation.
\end{itemize}

Suggestions were not accepted solely because they were generated by an AI
system. They were checked against source documentation, governing equations,
or executable evidence. Several examples demonstrate the necessity of this
critical process:
\begin{enumerate}
  \item the initial positive-weight parameterization produced weights near
        \(\operatorname{softplus}(0)\), causing excessive initial losses; the
        initialization was corrected and the model retrained;
  \item raw constitutive boundary-traction integration produced a persistent
        Hill--Mandel mismatch of \(3.859\times10^{-6}\); local reproduction
        showed that the correct discrete quantity required assembled particle
        reaction work, reducing the mismatch to approximately \(10^{-13}\);
  \item an integration-point diagnostic returned zero deformation gradients
        although the mixed solve was correct; direct coefficient evaluation
        isolated and corrected the sampling error.
\end{enumerate}

These episodes are reported because they distinguish assistance from
authority: the accepted implementation is defined by the equations and
verification evidence, not by chatbot output.
